\section{Shared Experimental Setup}
\label{sec:setup}

\subsection{Models and Data}
The documented model family is Qwen3-8B with the Red Hat EAGLE3 drafter. The
current suite compares the original drafter, a trained ordinary LoRA, and
generator snapshots labeled 2,000 and 4,000 steps. Generator data size is
reported as 24K examples. Exact training progress must ultimately come from
checkpoint metadata; labels are retained without treating them as verification.

Nemotron math, code, STEM, and chat are designated in-distribution evaluation
categories. Red Hat HumanEval, math reasoning, and summarization are designated
out-of-distribution categories in this comparison. These labels describe the
chosen evaluation sources; they do not by themselves establish absence of all
semantic overlap with training. Validation rows are excluded from optimization
by the documented trajectory membership contract.

The full-data workflow schedules six training subsets per domain, totaling
24K examples, and supports ten planned epochs. It saves Hub snapshots every
1,000 optimizer updates and at completion, with context-4 validation every 500
updates in the current recipe. The schedule is not evidence that ten epochs
have completed. This report includes only the supplied checkpoint observations.

\subsection{Metrics and Measurement Boundaries}
We keep the following quantities distinct:
\begin{itemize}[leftmargin=*]
  \item \textbf{Training and validation losses:} optimization/evaluation values
  under a specified conditioning and microbatch policy.
  \item \textbf{Offline agreement:} drafter/target argmax agreement in offline
  evaluation, not measured speculative acceptance.
  \item \textbf{Serving acceptance:} native speculative-decoding counters during
  actual generation; absent from the supplied current suite summary.
  \item \textbf{Throughput:} aggregate generated output tokens divided by
  measured wall seconds, including prefill and serving overhead.
  \item \textbf{Latency and systems timings:} request latency, first-token time,
  and stage wall time, when their artifacts are available.
\end{itemize}

\subsection{Recipe Versus Verified Run Metadata}
The independent serving implementation fixes BF16, eager execution, the V1
runner, and a pinned vLLM 0.31.0 integration. Suite defaults specify a seed,
sample limit, output budget, and warmup count. These defaults must not be
substituted for missing metadata from the supplied throughput summary.
For that summary, hardware, concurrency, sample counts, seed, workload hashes,
and exact checkpoint identities remain unverified. The cache pilot separately
reports an A100 80 GB. Neither source establishes repeated-run uncertainty.
